\section{What is proved and what is open}\label{sec:status}

\subsection{Proved}

Inside the model of Section~\ref{sec:model}, the paper determines the least cost of soundness. For finitely many candidates and exponential checks it is the entropy law, exact up to an additive term that is sharp (Theorem~\ref{thm:A}) and exact without remainder when the soundness is the reciprocal of an integer and the budget a power of it, where the optimal schedule is a Huffman code (Theorem~\ref{thm:C}). For checks of polynomial order the least cost is given in closed form by a R\'enyi entropy (Theorem~\ref{thm:D}), and for every cost function an infinite sequence of candidates admits a sound schedule of finite cost exactly when one series built from the concave conjugate of the cost converges (Theorem~\ref{thm:dual}). The same price, rescaled by the error rate of the generator, governs chains of improvements (Theorem~\ref{thm:B}), and checking only the outcome of a chain of $n$ steps is exponentially more expensive than checking every step (Theorem~\ref{thm:gating}). An infinite chain reaches a singularity soundly exactly when the costs of its checks have finite entropy in addition to a finite sum, which fails for a nonempty range of cost laws (Theorem~\ref{thm:E}, Corollary~\ref{cor:gap}); a branching takeoff with a fixed false acceptance rate keeps infinitely many errors on survival, and a sound one needs a number of runs that grows linearly with the generation (Theorem~\ref{thm:F}). In the Osgood model of growth every singular takeoff can be made sound with exponential checks, at a price of order $1/(\alpha-1)^{2}$ near the critical exponent, while polynomial checks move the threshold of a sound takeoff from $\alpha=1$ to $\alpha=1+s$ (Proposition~\ref{prop:osgood}).

The finite inequalities behind these statements are proved in Lean~4 (Appendix~\ref{app:checks}). The probabilistic limit theorems, Wald's identity, the Kesten--Stigum theorem, L\'evy's conditional Borel--Cantelli lemma and the extinction theorem of Galton--Watson processes, are classical and are proved or cited in Appendix~\ref{app:derivations}; the asymptotics of $\zeta$ near $1$ are proved in the text.

\subsection{Open}

Three questions are left open, and each is a precise mathematical question.

\emph{Adaptive protocols.} Remark~\ref{rem:adaptive} shows that a protocol which looks at the outcomes of the checks can avoid the entropy term by halting at the first rejection, at the price of losing the correct work after it. Let a protocol be \emph{conservative} if, in every world, it keeps every correct candidate with probability one. Is the least worst-case expected cost of a conservative $\theta$-sound protocol equal to $\Vs(\theta)$ up to an additive $O(A)$? An adaptive protocol may let the number of runs of later checks depend on the results of earlier ones; since a correct candidate never fails, a conservative protocol learns nothing from the runs of a candidate that passes, and the question is whether the information it gains from rejections can be converted into a saving that is not $O(A)$.

\emph{Two-sided checks.} For checks that err on both sides Proposition~\ref{prop:twosided} gives the entropy law with the rate $D=\mathrm{KL}(\mathrm{Ber}(\frac12)\,\|\,\mathrm{Ber}(q))$ in the upper bound and $\log(1/q)$ in the lower bound. The exact rate of a sequential test with two error probabilities is Wald's \cite{Wal45}, and the question is whether the optimal schedule for many two-sided checks attains the entropy law with the Chernoff rate of the pair of hypotheses, in both directions.

\emph{Correlated runs.} The model assumes that the runs of a check are independent. A real verifier, a learned reward model or a test suite run on the same machine, has correlated errors, and repeated runs of a deterministic check gain nothing. The entropy law holds for every family of checks whose false acceptance probability after $r$ runs is at most $\eps^{r}$, whatever the mechanism; the open question is which checks used in practice have this property, with which $\eps$, and what replaces the law for those that do not.

\subsection{What is not claimed}

The paper does not claim that any existing system is sound or that an intelligence explosion will or will not occur. Its theorems are statements about a model with independent runs, perfect completeness and schedules fixed in advance, and their value for a real system depends on how well that system satisfies them. What is claimed is that inside the model the cost of soundness is computable, that it is an entropy, and that it divides the singular processes into those that can be made sound and those that cannot.
